\documentclass[a4paper, oneside]{article}
\usepackage{graphicx}
\usepackage{amsthm}
\usepackage{amsmath}
\usepackage{amssymb}
\usepackage[a4paper,
            bindingoffset=0.2in,
            left=2cm,
            right=2cm,
            top=2cm,
            bottom=2cm,
            footskip=.25in]{geometry}
\usepackage[italian]{babel}
\usepackage{pgfplots}
\usepackage{tabularx}
\usepackage{tikz}
\usepackage{wrapfig}
\usepackage{color}
\usepackage[d]{esvect}
\definecolor{page}{rgb}{0.129,0.157,0.212}
\pagecolor{page}
\color{white}
\graphicspath{ {./images/} }
\usetikzlibrary{shapes.geometric}
\usetikzlibrary{datavisualization}
\usetikzlibrary{datavisualization.formats.functions}
\usetikzlibrary{patterns}
\pgfplotsset{width=10cm,compat=1.18}

\title{Appunti di Fisica II}
\author{Tommaso Miliani}
\date{12-03-26}

\begin{document}
\newtheoremstyle{theoremEnv}
                {}          % Space above
                {}          % Space below
                {\slshape}  % Body font
                {}          % Indent amount
                {\bfseries} % Head font
                {.}         % Punctuation after head
                {\newline}  % Space after theorem head
                {}          % Theorem head spec
\theoremstyle{theoremEnv}

\newtheorem{definition}{Definizione}[section]
\newtheorem{theorem}{Teorema}[section]
\newtheorem{lemma}{Proposizione}[section]
\newtheorem{observation}{Osservazione}[section]
\newtheorem{corollary}{Corollario}[theorem]
\newtheorem{example}{Esempio}[section]
\newtheorem{remark}{Enunciato}[section]

\maketitle

\section{Condensatore con carica}
Data una lastra metallica di spessore $D$ tra le due piastre
del condensatore, in modo tale che sia distante $a$ dalla prima 
piastra e $b$ dalla seconda, dunque $a + b = D$. 
La piastra centrale avrà, per induzione, $-Q_0$ sulla parte sinistra
(in quanto la prima piastra ha $Q_0$) e $Q_0$ sulla destra
(in quanto la piastra a destra ha $-Q_0$). Dato che la piastra rimane neutra, 
la somma delle sue cariche interne è zero, inoltre essa con le due
piastre costituisce un sistema di due condensatori in serie. 
\begin{gather*}
    \frac{1}{C_1} = \frac{a}{\epsilon_0 S} + \frac{b}{\epsilon_0 S} = \frac{a + b}{\epsilon_0 S} = \frac{D}{\epsilon_0 S} \ \Longrightarrow \ C_1 = \frac{\epsilon_0 S}{D} = 2C_0
\end{gather*}
Il processo di carica di un condensatore consiste nel prendere singolarmente le cariche da una piastra e spostandole
nell'altra piastra in modo tale che da una parte si abbia $-Q_0$ e $Q_0$ fino a 
che la differenza di potenziale tra le due piastre è 
\begin{gather*}
    V_0 \frac{Q_0}{C_0}
\end{gather*}
Il lavoro compiuto nel portare un $dq$ dalla piastra carica con $q$ a quella $-q$, si 
compirà un lavoro 
\begin{gather*}
    \mathcal{L} = \int_{\gamma} dq\vv{E} \cdot \vv{dl} = dq \int_{\gamma} \vv{E} \cdot \vv{dl} = -dq \Delta V  = dq(V(D) - V(0)) = dq(-\Delta V)
\end{gather*}
Dove $\Delta V$ è la differenza di potenziale tra le armature quando è presente
una carica $q$. Si osserva che il campo compie un lavoro negativo (poiché 
si oppone al trasporto della carica $dq$ da un armatura all'altra), mentre il generatore
compie lavoro positivo per portare le cariche da un armatura all'altra. Se si ponesse $-dq$:
\begin{gather*}
    \mathcal{L}_{E} = -\int_{\gamma} dq\vv{E} \cdot \vv{dl} = -dq \int_{\gamma} \vv{E} \cdot \vv{dl} = dq \Delta V  = dq(V(D) - V(0)) = dq(-\Delta V)
\end{gather*}
Mentre il lavoro del generatore
\begin{gather*}
    d\mathcal{L}_{G} = - d\mathcal{L}_{E} = dq(V(0)) = dqV_0
\end{gather*}
Quando si mette a massa la piastra destra a massa  $V(D) = 0$. Adesso 
il lavoro meccanico complessivo diventa
\begin{gather*}
    \mathcal{L}_{G} = \int_{0}^{Q_0}  
\end{gather*} 
Più cariche si vogliono mettere sull'armatura destra e maggiore 
sarà il lavoro richiesto al generatore per portare le cariche successive, 
dunque $V(0)$ cresce nel tempo. Allora il lavoro si può esprimere in funzione 
della carica già presente $q$:
\begin{gather*}
    dqV(0) = dq\frac{q}{\mathcal{C}_0 }
\end{gather*}
Allora 
\begin{gather*}
    \mathcal{L}_{G} = \int_{0}^{Q_0} dq\frac{q}{\mathcal{C}_0} = \frac{Q_0^{2}}{2\mathcal{C}_0}  = \frac{Q_0V_0}{2} = \mathcal{C}_0 \frac{V_0^{2}}{2}
\end{gather*}
Questa energia viene accumulata nel campo elettrico (che prima non c'era).
ispetto al 
caso in cui il condensatore non era carico, il generatore ha  compiuto un lavoro 
e ha immagazzinato quell'energia in un campo elettrico secondo la seguente:
\begin{gather*}
    \frac{Q_0^{2}}{2\mathcal{C}_0} = \frac{\epsilon E^{2}}{2}S(2D)
\end{gather*}
Se si inserisse la lastra metallica tra le due armature, allora
\begin{align*}
    U_1 &= \frac{Q_0^{2}}{2\mathcal{C}_1} = \frac{U_0}{2}  \qquad &Q_0 = \text{const}\\
    U_1 &= \frac{\mathcal{C}_1 V_0^{2}}{2} = 2U_0 \qquad &V = \text{const}
\end{align*}
Quale delle due espressioni è vera? Sono entrambe vere ma dipendono dal fatto che, la prima 
è vera se si tiene stabile la carica $Q_0$, mentre la seconda vale se $V$ rimane costante. 
Se le due armature non sono conesse a niente, le cariche rimangono costanti 
ed il lavoro di inserimento della lastra è tale per cui l'energia dimezzi: questo perché 
dentro la lastra non c'è campo elettrico e, dato che la dimensione della lastra ($D$) è uguale alla distanza
tra l'armatura 1 e la lastra + l'armatura 2 e la lastra, allora si è dimezzato il volume 
in cui è presente il campo. Il campo elettrico ha dunque speso energia nell'inserimento 
della lastra (cedendola a quest'ultima sottoforma di energia cinetica e quindi compiendo lavoro
positivo). I contributi vicino al bordo non possonon essere bilanciati, allora le linee di campo 
ai bordi del condensatore non sono più orizzontali, altrimenti la lastra non 
potrebbe essere tirata dentro dal campo elettrico. Sulle superfici 
della lastra agiscono due forze concordi verso il basso ma opposte sulla direzione 
orizzontale, questa forza è responsabile del dimezzamento dell'energia potenziale:
\begin{gather*}
    \mathcal{L}_{E} = \frac{U_0}{2}
\end{gather*}
IL lavoro di portare una carica da un armatura all'altra lo si fa mediante un \textbf{generatore}: ossia un 
dispositivo che riesce a fare un lavoro contrario alle forze del campo facendo 
scorrere una corrente che muove gli elettroni con una forza motrice (non elettrica). 
Invece nel caso in cui l'energia cresce, il generatore sposta le cariche sulle armature 
in modo tale che non scenda il potenziale nel condensatore. Riassumendo
\begin{itemize}
    \item Carica costante:
    \begin{gather*}
        U_{in}^{Q} = U_0^{Q} \qquad U_{fin}^{Q} = \frac{U_0}{2} 
    \end{gather*}
    Allora
    \begin{gather*}
        \Delta U^{Q} = - \frac{U_0}{2} = - \mathcal{L}_E
    \end{gather*}
    In questo caso $V = \frac{V_0}{2}$.
    \item Potenziale costante: 
    \begin{gather*}
        U_{in}^{V} = U_0 \qquad U_{fin}^{V} = 2U_0
    \end{gather*}
    Allora
    \begin{gather*}
        \Delta U^{V} = U_0 =  \mathcal{L}_G - \mathcal{L}_E
    \end{gather*}
    In questo caso $Q = \frac{Q_0}{2}$. 
\end{itemize}
Si può dunque definire seguenti procedimenti:
\begin{enumerate}
    \item Inserire la lastra nel condensatore, in questo modo 
    diminuisce l'energia potenziale ma la carica rimane costante:
    $Q_0 = Q$, ma $V = \frac{V_0}{2}$. 
    \item Si attacca il generatore in modo tale da riportare $V$ a $V_0$, 
    dunque $Q = 2Q_0$. Il lavoro compiuto dal generatore è ovviamente positivo 
    poiché sta aumentando l'energia potenziale del condensatore.
\end{enumerate}
Il lavoro compiuto dal generatore è dunque 
\begin{gather*}
    \mathcal{L}_G = \int_{Q_0}^{2Q_0} dq V(q)  \qquad V(q) = \frac{q}{2\mathcal{C}_0} = \frac{q}{\mathcal{C}_1}
\end{gather*}
Dove il potenziale $V(q)$ sta salendo man mano che le cariche spostate 
aumentano, allora l'integrale diventa
\begin{gather*}
    \mathcal{L}_G  = \int_{Q_0}^{2Q_0} \frac{q}{\mathcal{C}_1} dq = \frac{3Q_0^{2}}{2\mathcal{C}_1} = \frac{3}{2}U_0 
\end{gather*}

\section{Dielettrici}
\subsection{Sviluppo multipolare dei dielettrici}
Considerato un volume $\tau$ molto grande all'interno del quale è presente
della carica con densità di carica $\rho$ (all'esterno non vi è carica),
si considera un sistema di assi centrato dentro $\tau$. Si vuole ora
determinare il campo elettrico a distanza $\vv{r}$. 
\begin{gather*}
    V(r) = \frac{1}{4\pi\epsilon_0} \int \frac{\rho(\vv{r'} )}{\left| \vv{r} - \vv{r'}   \right| } \ d\tau
\end{gather*} 
Dato che $d\tau \rho = dq$. Se si è abbastanza lontani dal volumetto dove sta la carica, 
la distribuzione precisa della carica non è importante per determinare il 
potenziale, ma è possibile eseguire uno sviluppo in serie del potenziale in modo che, 
via via che si sviluppa aumenta la precisione del potenziale. Posti 
allora nel caso in cui $\left| \vv{r}  \right|  >> \left| \vv{r'}  \right|$, si scrive
\begin{gather*}
    \frac{1}{\left| \vv{r} -\vv{r'}   \right| } = \frac{1}{\left((\vv{r} - \vv{r'}  ) \right)^{1/2}} = \left(r^{2} + r^{'2} - 2(\vv{r}\cdot \vv{r'}  )\right)^{\frac{1}{2}} = \frac{1}{r} \left(1 + \underset{x}{\underline{\frac{r'^{2}}{r^{2}} - 2\frac{(\vv{r} \cdot \vv{r'}  )}{r^{2}}}}\right)^{\frac{1}{2}}
\end{gather*}
Dunque si pul ottenere lo sviluppo in funzione dell'infinitesimo $x$:
\begin{gather*}
    \frac{1}{r}\left(1 - \frac{1}{2}x  + \frac{3}{8}x^{2}\right)
\end{gather*}
Tuttavia dentro $x$ c'è anche un infinitesimo del secondo ordine, allora per fermarsi al primo ordine devo 
fermarmi al secondo ordine di $\frac{r'}{r}$:
\begin{gather*}
    \frac{1}{\left| \vv{r} - \vv{r'}   \right| } = \frac{1}{r}\left(1 - \frac{1}{2} \frac{r^{'2}}{r^{2}} + \frac{(\vv{r} -\vv{r'}  )}{r^{2}} + \frac{3}{8} \left(4\frac{\left(\vv{r} - \vv{r'}  \right)^{2}}{r^{4}}\right) + o\left(\left(\frac{r'}{r}\right)^{3}\right)\right)
\end{gather*}
Allora il potenziale diventa
\begin{gather*}
    V(r) = \frac{1}{4\pi\epsilon_0} \int \rho(r') \frac{1}{\left| \vv{r} -\vv{r'}   \right| } d\tau \approx \frac{1}{4\pi \epsilon_0}\left(\int \frac{\rho(r')}{r} dr  + \int \frac{\rho(r')}{r^{3}} (\vv{r} \cdot \vv{r'}  ) \ d\tau + \frac{1}{2} \int \left(\frac{3(\vv{r}\cdot \vv{r'}  )^{2}}{r^{5}} - \frac{r'^{2}}{r^{3}}\right) \rho \ d\tau  \right) 
\end{gather*}
Questo è definito come \textbf{sviluppo multipolare} del potenziale:
\begin{gather*}
    V(r) = \frac{1}{4\pi\epsilon_0} \frac{1}{r}\int \rho(r') \ d\tau = \frac{Q}{4\pi\epsilon_0 r} 
\end{gather*}
Questo prende il nome di \textbf{monopolo}: da lontano tutte le distribuzioni di carica
appaiono come puntiformi data la distanza e dunque la carica $Q$ è il \textbf{monopolo elettrico}, 
fermandoci per un termine diverso dall'ordine uno poiché altrimenti è zero (???). Per l'ordine 
1 si prende il termine dello sviluppo:
\begin{gather*}
    V_1(r) = \frac{1}{4\pi\epsilon_0} \frac{1}{r^{3}} \int \rho(\vv{r'} ) \cdot \vv{r} \cdot \vv{r'} \ d\tau   = \frac{1}{4\pi\epsilon_0} \frac{\vv{r} }{r^{3}} \int \rho(\vv{r'} ) \vv{r'} \ d\tau  
\end{gather*}
Posso portare fuori il vettore $\vv{r}$ in quanto è sempre costante. L'integrale finale prende il nome di 
\textbf{dipolo elettrico} ($\vv{p}$ ): se la carica è distante questo vettore va dalla carica negativa a quella positiva 
e da lontano in un dipolo elettrico cambia l'integrale a seconda della distribuzione di cariche positive e negative. Se $Q = 0 $, 
si può dire che $\rho = \rho_+(\vv{r'} ) + \rho_- (\vv{r'} )$, questo vuol dire che
\begin{gather*}
    \int_{\tau} \rho_-(\vv{r'} ) \ d\tau = Q \qquad \int_{\tau} \rho_-(\vv{r'} ) \ d\tau = -Q  
\end{gather*}
Il dipolo elettrico è dunque importante quando la carica totale è nulla:
\begin{gather*}
    \vv{p} = \int \left(\rho_+(\vv{r'} ) + \rho_-(\vv{r'} )\right) \vv{r} \ d\tau   = q\vv{r_+} - q\vv{r_-} = q(\vv{r_+} - \vv{r_-}  )  
\end{gather*}
Quando il dipolo contiene una carica $+$ in $\vv{r_+}$ ed una carica $-$ in $\vv{r_-}$, 
allora il vettore dipolo corrisponde esattamente al vettore dalla carica meno a quella più. Dunque 
quando si è molto lontani, non importa sapere dove stanno i centri di carica più e quelli di carica meno,
ma interessa solo che ci sia una carica positiva e negativa. Si capisce che ora, rispetto alla distribuzione 
di carica, conta rispetto a quale direzione rispetto all'asse del dipolo ci si allontana.
Nel caso in cui sia presente un altro dipolo uguale e contrario rispetto a quello di partenza,
si deve sviluppare al secondo ordine:
\begin{gather*}
    V_2(\vv{r} ) = \frac{1}{4\pi\epsilon_0} \frac{1}{2r^{5}}\int\rho\left(3(\vv{r} \cdot \vv{r'}  )^{2} - r^{'2} r^{2}\right) d\tau
\end{gather*}
Si tira fuori la matrice adesso utilizzando gli indici:
\begin{gather*}
    \frac{1}{4\pi\epsilon_0} \frac{1}{2r^{5}} \int \rho \left(3(\vv{r} \cdot \vv{r'}  ) (\vv{r}\cdot \vv{r'}  ) - r^{'2} r^{2}\right) d\tau =  \frac{1}{4\pi\epsilon_0} \frac{1}{2r^{5}} \int \rho \left(3(r_ir_i')(r_jr_j) - r'^{2} r^{2}\right) d\tau \qquad r^{2} = \sum_{i, j = 1}^{n} r_ir_j \delta_{i, j}
\end{gather*}
Il motivo di $j$ è che devo moltiplicare ma non riutilizzare gli stessi indici, in questo modo esce fuori la matrice, 
inoltre si mette anche lo scalare in matrice:
\begin{gather*}
    V_2(r) = \frac{1}{4\pi\epsilon_0} \frac{1}{2r^{5}} r_i r_j \int \rho\left(3r_i'r_j' - r'^{2}\delta_{i,j}\right)  \ d\tau
\end{gather*}
L'integrale prende ora il nome di \textbf{tensore di quadrupolo}, definito dunque come
\begin{gather*}
    \mathcal{Q}_{i, j }\frac{1}{2} \int \rho (3r_i'r_j' - r'^{2}\delta_{i, j}) \ d\tau 
\end{gather*}
Si potrebbe osservare come in realtà esistano anche dipoli orizzontali (come
nel disegno), da qui il nome di \textbf{quadrupolo}. Allora si riassume lo sviluppo al secondo ordine come 
\begin{gather*}
    V_2(r) = \frac{1}{4\pi\epsilon_0} \frac{r_ir_j Q_{i, j}}{r^{5}}
\end{gather*}
Se questo fosse nullo, allora si otterrebbe l'ottopolo. Da qui si capisce il nome di sviluppo
multipolare. 

\subsection{I dielettrici}
Così come i metalli, in media le cariche si orientano o deformano fisicamente o sotto l'effetto di un campo 
elettrico in modo tale da ottenere un dipolo polarizzato (anche se di base non è polare).






\end{document}